\chapter{El Pipeline de Emergencia Cósmica}
La teoría se valida mediante la simulación de la cadena de complejidad, desde las fluctuaciones cuánticas primordiales hasta la aparición de civilizaciones tecnológicas.

\section{TCCU Galaxy Genesis Engine}
La formación de galaxias ocurre cuando el contraste de densidad $\delta = (\rho - \bar{\rho})/\bar{\rho}$ supera el umbral crítico $\delta_c$, generando un halo gravitacional. En el marco TCCU, estos halos son interpretados como \textbf{nodos de alta coherencia informacional}.

\section{TCCU Stellar Forge Engine}
Las estrellas emergen de nubes moleculares que superan la \textbf{Masa de Jeans}:
\begin{equation}
M_J = \left(\frac{5kT}{G\mu m_H}\right)^{3/2} \left(\frac{3}{4\pi\rho}\right)^{1/2}
\end{equation}
El simulador asigna masas siguiendo una ley de potencia (IMF): $\xi(M) \propto M^{-2.35}$.

\section{TCCU Life Emergence Engine}
La vida aparece en planetas dentro de la zona habitable ($r = \sqrt{L/L_\odot}$). El proceso se modela como una optimización de la complejidad química $C = \sum p_i \ln p_i$ y la aparición de replicadores mediante la ecuación logística:
\begin{equation}
\frac{dN}{dt} = rN \left(1 - \frac{N}{K}\right)
\end{equation}

\section{Sincronización Multiescalar}
La vida en la Tierra (emergida hace 3.7 Gyr) coincide con las primeras 4-5 órbitas galácticas del Sol. TCCU postula que el cruce de los brazos espirales actúa como un \textbf{modulador informacional} que estabiliza la química planetaria.
